\section{結果と考察}

% 追加: 主観評価(7段階)を横バーで可視化するマクロ。第1引数=平均値(0-7)。
% 全長 \likertfull を 7 段階に対応させ，値に比例した塗り(濃灰)+残り(淡灰)で表示する。
\newlength{\likertfull}\setlength{\likertfull}{24mm}
\newcommand{\likert}[1]{%
  \begingroup
  \setlength{\fboxsep}{0pt}%
  \colorbox[gray]{0.45}{\rule{0pt}{7pt}\hspace{\dimexpr #1\likertfull/7\relax}}%
  \colorbox[gray]{0.88}{\rule{0pt}{7pt}\hspace{\dimexpr \likertfull - #1\likertfull/7\relax}}%
  \endgroup}

\subsection{入力効率（KSPC）}
入力効率は提案手法（母音入力）の KSPC が $1.12$ であり，
ローマ字入力（訓令式，$1.95$），JIS かな入力・フリック入力（いずれも $1.34$），トグル入力（$3.06$）より少なく，
打鍵数の面で省入力が達成できている（表\ref{tab:input}）．

\begin{table}[t]
  \centering\small
  \caption{入力効率比較（KSPC は理想 KSPC＝誤りなし最小打鍵数 $\div$ 文字数．評価セット 1{,}378 文上で算出）}
  \label{tab:input}
  \setlength{\tabcolsep}{4pt}
  \begin{tabular}{p{3.4cm}cc}
    \toprule
    Method & Keys & KSPC \\
    \midrule
    提案手法（母音入力 + LoRA LLM）& 5 + \texttt{n} & 1.12 \\
    ローマ字入力（訓令式）& 26+ & 1.95 \\
    JIS かな入力 / フリック入力 & 12--50 & 1.34 \\
    トグル入力 & 12 & 3.06 \\
    \bottomrule
  \end{tabular}
\end{table}

\subsection{精度評価とアブレーション}
\label{subsec:acc}
文全体の完全一致精度は，基準条件 a（CTXあり・ビーム 8，提案手法）で
Acc@1 $=7.0\%$，Acc@3 $=11.8\%$，Acc@5 $=14.1\%$ と低く（表\ref{tab:abl}），
母音化による情報欠落の大きさを反映している．これは完全一致ではなく候補提示型支援としての位置づけの妥当性を示す．
LoRA を外した素のベースモデル（条件 e）は Acc@1 が 0\% であり，本タスクが LoRA による追加学習に強く依存することが分かる．
ビーム幅は 4（条件 c）より 8（基準 a），さらに 16（条件 d）の方が高精度であった．
CTX の有無（a と b）では CTX ありが一貫して上回り，CTX の付与が精度向上に寄与することが確認できた（本評価セットは CTX を持つ文を約 32\%＝483 文含む）．

\begin{table}[t]
  \centering\footnotesize
  \caption{精度評価・アブレーション（評価セット 1{,}378 文上の完全一致 Acc@$k$．基準 a が提案手法）}
  \label{tab:abl}
  \setlength{\tabcolsep}{4pt}
  \begin{tabular}{@{}p{4.0cm}ccc@{}}
    \toprule
    Condition & Acc@1 & Acc@3 & Acc@5 \\
    \midrule
    a：CTXあり・ビーム 8（提案手法）& 7.0\% & 11.8\% & 14.1\% \\
    b：CTXなし・ビーム 8 & 5.4\% & 10.0\% & 12.2\% \\
    c：CTXあり・ビーム 4 & 7.0\% & 10.7\% & 11.8\% \\
    d：CTXあり・ビーム 16 & 7.5\% & 13.1\% & 15.2\% \\
    e：LoRA なし（素のベース）& 0.0\% & 0.0\% & 0.0\% \\
    \bottomrule
  \end{tabular}
\end{table}

\subsection{実ユーザ実験結果}
\label{subsec:userresult}
参加者 5 名による各方式・各 Level の試行（正答試行のみ）について，
総入力時間 $T_{\mathrm{total}}$・初動時間 $T_{\mathrm{init}}$・打鍵数の
参加者平均を表\ref{tab:userobj}に，母音入力の主観評価（7 段階）の
参加者平均を表\ref{tab:usersubj}に示す。

\begin{table}[t]
  \centering\small
  \caption{客観指標の参加者平均（$N=5$，正答試行のみ）。母音入力＋推論は打鍵がほぼ 1 で総入力時間が最短。}
  \label{tab:userobj}
  \setlength{\tabcolsep}{4pt}
  \begin{tabular}{p{3.2cm}ccc}
    \toprule
    Method & $T_{\mathrm{total}}$[ms] & $T_{\mathrm{init}}$[ms] & 打鍵数 \\
    \midrule
    ひらがな入力& 6495 & 1224 & 14.1 \\
    母音入力& 6040 & 1539 & 8.2 \\
    標準IME& 6675 & 1100 & 14.9 \\
    母音入力＋LLM推論& 3489 & -- & 1.2 \\
    \bottomrule
  \end{tabular}
\end{table}

\begin{table}[t]
  \centering\footnotesize
  \caption{母音入力の主観評価の参加者平均（$N=5$，7 段階：1=全くそう思わない～7=非常にそう思う）。バーは値に比例（最大 7）。母音変換負荷は高く，候補性能は高く評価された。}
  \label{tab:usersubj}
  \setlength{\tabcolsep}{4pt}
  \begin{tabular}{p{3.8cm}cl}
    \toprule
    項目 & 平均 & \\
    \midrule
    母音変換負荷（変換が難しい）& 4.6 & \likert{4.6} \\
    記憶保持負荷（覚えて入力が難しい）& 3.2 & \likert{3.2} \\
    操作負荷（操作が難しい）& 2.4 & \likert{2.4} \\
    混乱・ストレス & 5.4 & \likert{5.4} \\
    疲労 & 4.6 & \likert{4.6} \\
    \midrule
    候補が出た & 5.4 & \likert{5.4} \\
    候補が上位に出た & 5.2 & \likert{5.2} \\
    候補への信頼 & 4.6 & \likert{4.6} \\
    候補の有用感 & 4.4 & \likert{4.4} \\
    利用意向 & 3.2 & \likert{3.2} \\
    \bottomrule
  \end{tabular}
\end{table}

\subsection{考察}
\subsubsection{推論なしの比較}
まず，推論を介さない 2 方式を比べる。母音入力（推論なし）は
ひらがな入力 に対して打鍵数を $14.1$ から $8.2$ へ確実に削減する一方，
総入力時間は $6495$ ms 対 $6040$ ms とほぼ同等にとどまった。
すなわち母音化によって削減された物理的な打鍵コストは，
文を母音列へ変換し保持する認知コストへと置き換わっており，
省打鍵がそのまま時間短縮には結びつかない。
実際，母音入力（推論なし）は初動時間が長く（$1539$ ms 対 $1224$ ms），
入力前の音韻変換に時間を要していることがうかがえる。

\subsubsection{推論ありの比較}
次に，推論ありの 2 方式を比べる。母音入力＋推論は標準 IME の
$6675$ ms に対し $3489$ ms と約半分の総入力時間で入力を完了し，
打鍵数も平均 $1.2$ とほぼ「母音列を一度打って候補を選ぶ」操作に集約された。
特に IME での入力に時間を要する参加者ほど短縮効果が大きく，
LLM 推論が母音入力の変換・確定コストを肩代わりすることで，所要時間が半減したことから
提案手法が IME として実用的な入力速度に達することが示された。

\subsubsection{ボトルネックの所在}
ただし，母音入力＋推論の初動時間は本実験の計測方式では正しく取得できない。
他方式は打鍵がそのまま確定文字となるため初動時間を音韻変換時間として分離できるが，
母音入力＋推論では打鍵・候補生成・選択・確定が連続し，最初の打鍵までの値が候補を待つ時間を含むため同じ意味で比較できない。
このため初動時間は値を示さず（表\ref{tab:userobj} の ``--''），以降の議論には用いない。
代わりに，母音入力に伴う負荷は推論なしの母音入力から論じる。
この傾向は主観評価とも整合し，母音変換負荷（$4.6$）は高く評価された。
客観・主観の双方から，母音入力のボトルネックは
音韻変換（母音変換）にあることが分かった。
一方で操作負荷は低く（$2.4$），候補が出る・上位に出る・信頼できるといった
候補性能は高く評価された（$5.4$／$5.2$／$4.6$）。
これは提案手法の中核である LLM 候補提示が肯定的に受容されたことを意味する。
他方で混乱・ストレス（$5.4$）や疲労（$4.6$）は高く，
利用意向は中程度（$3.2$）にとどまった。
以上から，提案手法は IME として実用的な速度を達成しつつも，
残る課題は入力操作ではなく母音変換の負担軽減にあると結論づけられる。
なお参加者間の個人差は大きく，本結果は予備的な知見である。
